\section{Introduction}
\label{sec:introduction}

La fabrication additive par dépôt de fil fondu (\emph{Fused Deposition Modeling}, FDM) permet de concevoir des composants complexes intégrant des cavités fonctionnelles, des canaux internes et des treillis alvéolaires allégés (\emph{infill}) inaccessibles aux procédés conventionnels d'usinage. Toutefois, l'intégrité géométrique et mécanique de ces architectures internes est fréquemment compromise par des défauts d'impression inhérents au procédé~: sous-extrusions locales, bouchages partiels de buse, délaminations entre strates et variations dimensionnelles induites par le retrait thermique de solidification.

Les techniques usuelles de métrologie de surface (profilométrie optique, microscopie confocale, machines à mesurer tridimensionnelles à palpeur mécanique) sont incapables de sonder l'intérieur d'une pièce opaque sans sectionnement destructif. De même, les méthodes ultrasonores souffrent d'une atténuation acoustique très élevée et d'échos parasites multiples aux interfaces polymère-air. La tomodensitométrie à rayons X (micro-CT) constitue ainsi l'unique méthode non destructive offrant une cartographie tridimensionnelle quantitative de la matière et du vide à l'échelle sub-millimétrique~\cite{kruth2011computed, carmignato2014industrial}.

L'objectif central de ce laboratoire est d'évaluer quantitativement les performances métrologiques d'un tomographe à rayons X de laboratoire pour la détection et la caractérisation de défauts internes de dimensions étalonnées dans des pièces en thermoplastique. Le projet s'articule en deux étapes expérimentales distinctes~:
\begin{enumerate}
    \item \textbf{Séance 1~: Mise au point méthodologique et optimisation tomographique}~: Utiliser une pièce témoin en poly(acide lactique) (PLA) imprimée avec une buse de $0{,}4~\text{mm}$ (modèle biomimétique de cheval) afin de résoudre l'ambiguïté de l'axe de rotation du banc d'imagerie, de calibrer avec précision le centre de rotation (\emph{Center of Rotation}, COR) à l'échelle sous-pixel, d'éliminer les artéfacts en anneaux du capteur CMOS et de valider la résolution des parois fines.
    \item \textbf{Séance 2 (Partie principale)~: Contrôle non destructif et métrologie sur cubes étalons}~: Scanner des cubes de référence en PLA comportant des défauts géométriques internes étalonnés (canaux cylindriques et encoches de diamètres connus). Les mesures dimensionnelles obtenues par tomodensitométrie (largeur à mi-hauteur FWHM et analyse de profils) sont comparées directement aux modèles numériques CAO/STL afin de quantifier l'exactitude de mesure, la répétabilité, le rapport contraste sur bruit (CNR) et la limite de détection en taille de défaut.
\end{enumerate}

Le présent rapport est structuré de la façon suivante. La section~\ref{sec:theorie} rappelle brièvement les lois physiques de l'atténuation X (Beer-Lambert), le principe de la rétroprojection filtrée (FBP) et la méthode de propagation des incertitudes de mesure. La section~\ref{sec:methodologie} détaille le banc tomographique Leybold Didactic et la chaîne de calcul numérique. La section~\ref{sec:resultats_cheval} résume les acquis de la phase d'optimisation sur le spécimen en PLA (séance 1). La section~\ref{sec:metrologie_cubes}, c\oe{}ur de l'analyse, présente les résultats métrologiques obtenus sur les cubes à défauts calibrés (séance 2). Enfin, la section~\ref{sec:discussion} analyse les causes physiques d'incertitude et la section~\ref{sec:conclusion} synthétise la capabilité métrologique du tomographe pour le contrôle qualité industriel.
